% Created 2026-11-17 ter 16:54
% Intended LaTeX compiler: lualatex
\documentclass[11pt,a4paper]{scrartcl}
\usepackage{estudo}

\usepackage{graphicx}
\usepackage{longtable}
\usepackage{wrapfig}
\usepackage{capt-of}
\usepackage[normalem]{ulem}
\usepackage{amsmath}
\usepackage{amssymb}
\author{Nicolas}
\date{2026-11-17}
\title{Métricas}
\begin{document}

\maketitle
\tableofcontents

\section{Definição}
\label{sec:orgeed2053}
Uma métrica \(m\) é uma função quantitativa que busca \emph{medir} determinado \emph{atributo} de uma \textbf{entidade}, \textbf{processo}, \textbf{artefato} ou \textbf{sistema}:

$$
m : E \to V
$$

onde \(E\) é o \textbf{domínio das entidades} e \(V\) é o \textbf{domínio dos valores}. Normalmente, \(V\) assume valores inteiros ou reais, \emph{i.e.}, \(V = \mathbf{R}\) ou \(V = \mathbf{N}\).

\begin{exemplo}
Em \textbf{Ciência da Computação}, algumas métricas comuns são \emph{complexidade ciclomática}, \emph{tempo de execução}, \emph{taxa de erro} e \emph{cobertura de código} (dentre muitas outras).
\label{métricas--exemplos-em-computer-science}
\end{exemplo}

Nos exemplos acima, e na definição apresentada, estamos aumentando a semântica de \emph{métrica} para além das definições em \textbf{espaços métricos}:

\begin{teorema}[title={Observação}]
O \emph{conceito de métrica} quando utilizado fora do domínio dos \textbf{espaços métricos} não necessita, obrigatoriamente, satisfazer todos os axiomas destes espaços.
\end{teorema}
\section{Veja também}
\label{sec:org8e99c4d}
\begin{itemize}
\item \texttt{[Métrica em um Conjunto]} (ou \emph{métrica espacial}, para a definição mais rigorosa em Matemática)
\item \texttt{[Métricas de Distância]}  (métricas para as quais os axiomas dos espaços métricos são necessários; utilizadas em Ciência da Computação: funções de distância que satisfazem os axiomas métricos, como aquelas utilizadas em \emph{nearest-neighbor search}, \emph{clustering}, relações de similariedade, \emph{metric embeddings} e outros.
\end{itemize}
\end{document}
